\documentclass{article}
\usepackage[left=2cm, right=2cm, top=2cm, bottom=2cm]{geometry}
\usepackage{graphicx}
\usepackage{amsmath}
\usepackage{amssymb}
\usepackage{amsthm}
\usepackage{fancyhdr}
\usepackage{verbatim}
\usepackage{listings}
\usepackage{xcolor}
\usepackage{pdfpages}
\usepackage{cancel}
\usepackage{physics}
\usepackage{esint}
\usepackage{braket}

\lstset{
    language=C++,
    basicstyle=\ttfamily\footnotesize,
    keywordstyle=\color{blue},
    commentstyle=\color{green},
    stringstyle=\color{red},
    numbers=left,
    numberstyle=\tiny\color{gray},
    frame=single,
    breaklines=true,
    captionpos=b,
}


\title{MTH 446: Lecture \# 3 (catching up)}
\author{Cliff Sun}

\newtheorem{theorem}{Theorem}[section]
\newtheorem{lemma}[theorem]{Lemma}
\newtheorem{definition}[theorem]{Definition}
\newtheorem{conjecture}[theorem]{Conjecture}
\newtheorem{proposition}[theorem]{Proposition}
\newtheorem{corollary}[theorem]{Corollary}
\newtheorem{one minute paper}[theorem]{One Minute Paper}

\pagestyle{fancy}
\lhead{\textbf{\thepage}\ \ \nouppercase{\rightmark}}
\chead{MTH 446: Lecture \# 3 (catching up)}
\rhead{Cliff Sun}

\begin{document}

\maketitle

\section*{Lecture Span}
\begin{enumerate}
    \item Properties of Argument
\end{enumerate}

Let $z \in \mathbb{C}$, $z \neq 0$, $\varphi \in \arg(z)$, and if $-\pi < \varphi \leq \pi$, then 
$\varphi$ is called the \underline{principal value} of $\arg(z)$ with notation 

\begin{equation}
    \varphi = \arg(z)
\end{equation}

Therefore, 

\begin{equation}
    \varphi = \arg(z) \iff \varphi \in \arg(z) \land -\pi  < \varphi \leq \pi
\end{equation}

As an example with
\begin{equation}
    \arg(i) = \{\pi/2 + 2\pi n, n \in \mathbb{Z}\}
\end{equation}

Then, use $\pi/2 \in \arg(z)$ as the \underline{principal value}. Thus, 

\begin{equation}
    \arg(i) = \pi/2
\end{equation}

Also, 

\begin{equation}
    \arg(-1) = \{\pi + 2\pi n\}
\end{equation}

Thus, the principal value of $-1$ is $\pi$. 

Also, $\arg(\overline{z}) = \arg(z)$ if $z$ is not a negative real number (i.e., $\varphi \neq \pi$). 

\begin{proof}
    Let $\varphi = \arg(z)$, then $-\pi < \varphi \leq \pi$. Since $z$ is not a negative real number, then $\varphi \neq 0$. Also, $-\varphi \in \arg(\overline{z})$ (by reflection), also 
    \begin{equation}
        -\pi < \varphi < \pi \implies -\pi < -\varphi < \pi
    \end{equation}
    Thus, we conclude that $-\varphi = \arg(\overline{z})$. 
    Therefore, 
    \begin{equation}
        -\varphi = -\arg(z) = \arg(\overline{z})
    \end{equation}
\end{proof}

\section*{Exponentiation of complex numbers}

Let $r = |z| = \sqrt{x^2 + y^2}$, then 

\begin{equation}
    z = r\left( \frac{x}{r} + i \frac{y}{r} \right)
\end{equation}

Note that 

\begin{equation}
    \left( \frac{x}{r} \right)^2 + \left( \frac{y}{r} \right)^2 = 1
\end{equation}

Also, there is a $\theta$ such that 

\begin{equation}
    \cos\theta = \frac{x}{r}, \sin\theta = \frac{y}{r}
\end{equation}

Thus, 

\begin{equation}
    z = r(\cos\theta + \sin\theta) = e^{i\theta}
\end{equation}

For $\theta \in \mathbb{R}$. In general, $z \in \mathbb{C}$ and $z \neq 0$, then 

\begin{equation}
    z = re^{i\theta}, \theta \in \arg(z)
\end{equation}

You can also multiply complex numbers together 

\begin{equation}
    e^{\i\theta_1} e^{i\theta_2} = e^{i(\theta_1 + \theta_2)}
\end{equation}

Also, $\theta_1 \in \arg(z_1), \theta_2 \in \arg(z_2)$, then $\theta_1 + \theta_2 \in \arg(z_1 + z_2)$. 

\section*{Important remark}

Note that $\theta_1 = \arg(z_1)$ and $\theta_2 = \arg(z_2)$, then it is not always the case that $\theta_1 + \theta_2 = \arg(z_1z_2)$

List of important results:

\begin{enumerate}
    \item $e^{i(\theta_1 + \theta_2)} = e^{i\theta_1}e^{i\theta_2}$
    \item $e^{i(\theta_1 - \theta_2)} = e^{i\theta_1}/e^{i\theta_2}$
    \item $e^{-i\theta} = 1/e^{i\theta}$
    \item $|e^{i\theta}| = 1$
\end{enumerate}


\end{document}